\documentclass[11pt,a4paper]{article}
\usepackage[T1]{fontenc}
\usepackage[utf8]{inputenc}
\usepackage{lmodern}
\usepackage[margin=23mm]{geometry}
\usepackage{amsmath,amssymb,bm}
\usepackage{booktabs,longtable,tabularx,array}
\usepackage{graphicx,pdfpages}
\usepackage{xcolor,listings,enumitem}
\usepackage{microtype}
\usepackage[unicode,colorlinks=true,linkcolor=blue!55!black,
 citecolor=blue!55!black,urlcolor=blue!55!black]{hyperref}
\usepackage{bookmark}
\hypersetup{pdftitle={Single-top production with local recoil: physics, implementation and results},
 pdfsubject={MC@NLO, FKS, Pythia 8.318, stable-top t-channel production}}
\setlength{\parskip}{4pt}
\setlength{\parindent}{0pt}
\setlength{\emergencystretch}{2em}
\setcounter{tocdepth}{2}
\setlist{nosep,leftmargin=*}
\lstset{basicstyle=\ttfamily\small,breaklines=true,columns=fullflexible,
 keepspaces=true,showstringspaces=false,frame=single,framerule=.3pt,
 rulecolor=\color{black!25},aboveskip=7pt,belowskip=7pt}
\newcommand{\dd}{\mathrm{d}}
\newcommand{\as}{\alpha_{\mathrm{s}}}
\newcommand{\gs}{g_{\mathrm{s}}}
\newcommand{\CF}{C_{\mathrm{F}}}
\newcommand{\TR}{T_{\mathrm{R}}}
\newcommand{\pt}{p_{\perp}}
\newcommand{\zI}{z_{\mathrm{I}}}
\newcommand{\kt}{k_{\perp}}
\newcommand{\Born}{\Phi_{\mathrm{B}}}
\newcommand{\Real}{\Phi_{\mathrm{R}}}
\newcommand{\Dps}{\mathcal{D}_{\mathrm{PS}}}
\newcommand{\DG}{\mathcal{D}_{G}}
\newcommand{\code}[1]{\texttt{\detokenize{#1}}}
\newcommand{\file}[1]{\path{#1}}
\newcommand{\pb}{\,\mathrm{pb}}
\newcommand{\GeV}{\,\mathrm{GeV}}
\newcommand{\sampA}{default}
\newcommand{\sampB}{local}
\newcommand{\sampC}{local + collinear patch}
\newcommand{\sampD}{local + both patches}
\newcolumntype{L}[1]{>{\raggedright\arraybackslash}p{#1}}

\title{Single-top production with local recoil\\
\large Physics, FKS maps, MC@NLO counterterms, Pythia patches,\\
analysis, validation, and production plots}
\author{SingleTopLocal development record}
\date{24 October 2026}

\begin{document}
\maketitle
\begin{abstract}
This document records the development of local-recoil MC@NLO matching for
five-flavour, stable-top, $t$-channel single-top production in this workspace.
It gives the invariant recoil map and its measure, the combined initial--final
counterterms, the two singular-limit mismatches in stock Pythia 8.318, and the
implemented G-function and patched-shower alternatives. It also describes the
four process directories, their run settings, the 75-histogram analysis, the
shower-record truncation fix, and the numerical checks. The current four
million-event samples are compared quantitatively. The complete 73-panel
comparison atlas is embedded at the end, so this PDF contains the physics,
operational instructions, and plots in one place. Mathematical limits,
numerical evidence, and remaining accuracy qualifications are distinguished
throughout.
\end{abstract}
\tableofcontents
\clearpage

\section{Scope, outcome, and development decisions}\label{sec:scope}

The process is
\begin{equation}
 p p\longrightarrow t+j+X,
\end{equation}
with a stable on-shell top and NLO QCD corrections to electroweak $t$-channel
production. NLO subtraction uses FKS~\cite{FKS}, and shower matching uses
MC@NLO~\cite{MCatNLO,MG5}. The generation command is
\begin{lstlisting}
import model loop_sm-no_b_mass
define p = g u c d s b u~ c~ d~ s~ b~
define j = p
generate p p > t j $$ w+ w- [QCD]
\end{lstlisting}
The top sample is generated explicitly; a separate antitop production sample
is not included in the plotted results. Beam crossings and antiquark lines
are nevertheless included in the implementation tests. The process exclusion
in the command is part of the definition of this $t$-channel sample.

At Born level the colour flow contains two separate initial--final quark
connections: a light-quark line and the line connecting the incoming bottom
quark to the top. The colourless exchanged weak boson makes a local-recoil
study particularly natural. The implementation relies on this simple Born
structure and does not claim general-process support.

\begin{table}[htbp]
\centering\small
\begin{tabularx}{\linewidth}{@{}lXXX@{}}
\toprule
Process directory & Recoil / shower & Quark IF numerator & Local G replacement\\
\midrule
\code{pp_tj_default} & Ordinary MG5 MC@NLO / stock & Ordinary matching kernels & Ordinary global prescription\\
\code{pp_tj_local} & Local / stock & $1+\zI^2$ & Soft and outgoing-collinear\\
\code{pp_tj_local_patch} & Local / private one-patch build & $\zI^2+\zeta^2$ & Soft, with its collinear overlap\\
\code{pp_tj_local_2patch} & Local / private two-patch build & $(\zI^2+\zeta^2)r_t$ on top line & None of the special local replacements\\
\bottomrule
\end{tabularx}
\caption{The four implemented predictions. The common factor $\CF/(1-\zI)$
is suppressed in the numerator column. The two-patch light line uses
$\zI^2+\zeta^2$, without $r_t$. All samples retain ordinary FKS NLO subtraction.}
\label{tab:variants}
\end{table}

The work proceeded through the following substantive decisions.
\begin{enumerate}
\item Implement the shower's recoil map in FKS first, including forward and
inverse mappings, massive-top kinematics, phase-space Jacobians and finite
soft-endpoint terms. A correct map alone does not specify a matching kernel.
\item Treat the two ends of each Born IF colour line as one shower antenna,
then distribute its density over the two FKS orientations. Adding independent
ISR and FSR copies would double count the gluon-emission singularities.
\item Establish the stock-shower kernel before constructing subtraction.
Its outgoing massless-quark collinear coefficient and massive-top soft
current differ from exact QCD. These are kernel issues, not recoil-map issues.
\item Provide a stock-Pythia solution using local G replacements in MG5.
This restores finite event weights without silently changing Pythia.
\item Subsequently introduce two explicitly requested, opt-in alternatives:
a quark-collinear numerator patch, and then the additional top-soft patch.
Each has counterterms for the same modified shower rate.
\item Keep these alternatives in private Pythia installations. Shared stock
Pythia is never patched during a run, so no restoration race is introduced.
\item Use identical perturbative analysis and plotting definitions in all
four samples, and repair the independent HEPEVT history-capacity problem.
\end{enumerate}

An earlier development build modified a shower trial envelope, including
extra headroom. That implementation was withdrawn. Its results survive only
in the historical part of an audit JSON; they are not evidence for the
current stock-G or opt-in patch prescriptions. The current two kernel patches
are bounded suppressions and contain no such envelope enlargement.

The supported matching mode is ordinary MC@NLO for stable-top production,
without merging, resonance decays, QED corrections, or MC@NLO-Delta. Incoming
partons are massless in the analytic five-flavour calculation. The outgoing
top mass is retained. Pythia's own light/charm/bottom masses are not changed
to make the analytic approximation artificially exact.

\section{Notation and the common local recoil map}\label{sec:map}

\subsection{Invariant variables and the paper's notation}

Use the metric $p\cdot q=p^0q^0-\boldsymbol p\cdot\boldsymbol q$ and
future-directed physical incoming momenta. A tilde denotes the projected
Born state. For one selected colour connection, let
\begin{align}
 p_a^2&=0,& p_f^2&=m_f^2,& k^2&=0,\\
 \widetilde p_a^2&=0,& \widetilde p_f^2&=m_f^2.
\end{align}
Here $p_a$ is the real incoming quark, $p_f$ its real outgoing colour partner,
and $k$ the emitted gluon. The partner is either a massless light quark or
the top. Define
\begin{align}
 \zI&=1-\frac{p_f\cdot k}{p_a\cdot(p_f+k)},&
 w&=1-\zI,\label{eq:z}\\
 Q^2&=2p_a\cdot k,&
 D&=2\widetilde p_a\cdot\widetilde p_f,\label{eq:QD}\\
 \zeta&=\frac{p_a\cdot p_f}{p_a\cdot(p_f+k)}
       =1-\frac{\zI Q^2}{D},&
 \delta&=\frac{\zI Q^2}{D}=1-\zeta.\label{eq:zeta}
\end{align}
$D$ has mass dimension two; it is not a subtraction density. Densities are
denoted by calligraphic $\mathcal D$ below. The Bjorken relation is
$x_a^{\mathrm B}=\zI x_a^{\mathrm R}$.

The relation to Cabouat--Sj\"ostrand's notation~\cite{Cabouat} is
\begin{center}\small
\begin{tabularx}{\linewidth}{@{}lX@{}}
\toprule
Here & Paper / Pythia interpretation\\
\midrule
$\widetilde p_a$ & $p_b$, incoming before backward branching\\
$p_a$ & $p'_a$, incoming after backward branching\\
$\widetilde p_f$, $p_f$, $k$ & $p_f$, $p'_f$, $p'_c$ of the paper\\
$\zI$ & Backward fraction $z$; equal to $x$ in the massless DIS example\\
$\zeta$ & $z_1$ of that DIS example; the final-quark momentum fraction in its collinear limit\\
$D$ & $m_{\rm dip}^2-m_f^2$, with $m_{\rm dip}^2=(\widetilde p_a+\widetilde p_f)^2$\\
$wQ^2$ & Pythia evolution variable $p_{\perp\mathrm{evol}}^2$\\
$\kt^2$ & Physical transverse momentum, called \code{pT2corr} in the source\\
\bottomrule
\end{tabularx}
\end{center}
The identification with DIS $x,z_1$ is used only where the massless DIS
comparison applies; the invariant definitions above are used for the massive
top as well.

\subsection{Projection, mass preservation, and an explicit inverse}

Both FKS emitter orientations use the same real-to-Born projection:
\begin{equation}
 \boxed{\widetilde p_a=\zI p_a,\qquad
 \widetilde p_f=p_f+k-wp_a.}\label{eq:projection}
\end{equation}
All other physical momenta, including the other incoming parton, are fixed.
Equation~\eqref{eq:z} gives
\begin{equation}
 \widetilde p_f^2=m_f^2+2p_f\cdot k
             -2w p_a\cdot(p_f+k)=m_f^2,
\end{equation}
and momentum conservation follows from
$p_a-\widetilde p_a=p_f+k-\widetilde p_f$.
The top's recoil therefore changes its momentum while preserving its mass.

An invariant forward construction is also useful. Set
\begin{equation}
 n_f=\widetilde p_f-\frac{m_f^2}{D}\widetilde p_a,
 \qquad n_f^2=0,\qquad 2\widetilde p_a\cdot n_f=D.
\end{equation}
For a spacelike vector $k_\perp$ orthogonal to both basis vectors and with
$k_\perp^2=-\kt^2$, define
\begin{align}
 \kt^2&=wQ^2\zeta-m_f^2\delta^2,\label{eq:kt}\\
 A_k&=\frac{w}{\zI}(1-\delta)-\frac{m_f^2}{D}\delta,\\
 p_a&=\frac{\widetilde p_a}{\zI},&
 k&=A_k\widetilde p_a+\delta n_f+k_\perp,\\
 p_f&=\widetilde p_f+\frac{w}{\zI}\widetilde p_a-k.
\end{align}
The azimuth of $k_\perp$ fixes the remaining degree of freedom. These
relations imply $k^2=0$, $p_f^2=m_f^2$, and $2p_a\cdot k=Q^2$.
They express the same IF construction that the production generator
implements in its appropriate frames.

The physical interior satisfies
\begin{align}
 x_a^{\mathrm B}<\zI<1,\qquad
 0<\delta<\frac{wD}{wD+\zI m_f^2},\qquad
 \kt^2>0.\label{eq:support}
\end{align}
These are kinematic conditions. The fixed-order radiation map does not apply
a shower cutoff, SCALUP veto, or splitting probability to the real phase
space. In particular, shower termination must not remove the unresolved
region needed for NLO subtraction.

\subsection{Why the ISR generator samples \texorpdfstring{$w$}{w}, not the physical FKS energy}

The physical FKS variables retain their meaning in the real partonic centre
of mass: $\xi=2E_k/\sqrt{\hat s_{\mathrm R}}$ and $y$ is the cosine of the
gluon angle relative to the emitting beam. In general $\xi\ne w$. Solving for
the incoming rescaling at fixed $\xi$ can produce two solutions and a
turning-point Jacobian. Sampling $w$ directly covers the two solutions
continuously.

The implemented ISR measure can be stated explicitly. In the Born centre of
mass, write $e=\sqrt{\hat s_{\mathrm B}}/2$,
$f=\widetilde p_f$, $\eta=+1$ for beam 1 and $-1$ for beam 2,
$f_-=f^0-\eta f^3$, and $\lambda=1/(1-w)$. For the unit gluon direction
\begin{equation}
 n=(1,\sqrt{1-y^2}\cos\varphi,
        \sqrt{1-y^2}\sin\varphi,\eta y),
\end{equation}
define
\begin{align}
 A&=f_-(1+y)+2e(1-y),\\
 B_\varphi&=2(f^1\cos\varphi+f^2\sin\varphi)\sqrt{1-y^2},\\
 T_0&=2f\cdot n,\\
 T&=T_0+(\lambda-1)
             \left(A-\frac{B_\varphi}{\sqrt\lambda+1}\right),\\
 \xi&=\frac{2f_-\lambda w}{T}.
\end{align}
Here $B_\varphi$ is a geometric quantity, not the Born matrix element $B$.
With the PDFs and full-process flux supplied consistently by the caller,
the four-dimensional hadronic radiation measure is
\begin{equation}
 \dd\mu_{\mathrm{rad}}=
 \frac{\hat s_{\mathrm R}}{(4\pi)^3}\,
 \xi\,\frac{2f_-\lambda^2}{T}\,
 \dd w\,\dd y\,\dd\varphi.\label{eq:isrmeasure}
\end{equation}
The generator factors out the explicit $\xi$, as in the other FKS maps.
The real prefactor uses the physical $\xi$; the soft counterevent represents
the coefficient of $\dd w/w$ and uses $w$ instead. At the incoming-collinear
endpoint $y\to1$, $\xi=w$, and the ordinary ISR splitting fraction, PDF
remnants, and factorization scheme are retained.

Small scalar products are evaluated with rationalized masses and squared
direction differences. For example,
\begin{equation}
 f\cdot k=k^0\left[
 \frac{m_f^2}{f^0+|\boldsymbol f|}
 +\frac{|\boldsymbol f|}{2}
 \left|\frac{\boldsymbol f}{|\boldsymbol f|}
       -\frac{\boldsymbol k}{k^0}\right|^2\right].
\end{equation}
Using a model mass instead of subtracting $E^2-|\boldsymbol p|^2$ also avoids
inventing a small mass for a nearly beam-collinear massless partner.

\subsection{The finite soft-endpoint conversion}

For $w\to0$,
\begin{equation}
 \xi=a w+O(w^2),\qquad a=\frac{2f_-}{T_0}.
\end{equation}
In $4-2\epsilon$ dimensions the energy measure and reduced soft matrix
element bring a factor $a^{-2\epsilon}$. With an energy cutoff $h_s$,
\begin{equation}
 -\frac{(a h_s)^{-2\epsilon}}{2\epsilon}
 =-\frac1{2\epsilon}+\log h_s+\log a+O(\epsilon).
\end{equation}
The finite endpoint must therefore contain $\log a$. The implementation
uses \code{soft_cut=xicut/a} in the finite soft prefactor. The numerical
soft subtraction cutoff acts on $w$. Since $a\to1$ in the incoming-collinear
limit, the conversion is collinearly integrable and does not require a new
integrated pole operator or PDF counterterm.

\subsection{A common soft point for the combined S-event subtraction}

Different native FKS histories have different radiation frames, energy upper
bounds, and soft-energy slopes. Summing their subtraction terms at unrelated
soft gluon directions or slopes would leave a finite mismatch even if each
individual endpoint looked correct.

For local MC@NLO S events, the implementation uses one soft direction in the
Born centre of mass and converts it to each native frame, including the
aberration Jacobian. It also makes the physical soft energy agree. If a
native chart has relative soft slope $c>0$, the energy sampling variables
are related by the bijection
\begin{equation}
 u=r_E^2,\qquad
 v=\frac{u}{c(1-u)+u},\qquad
 r_E^{\mathrm{native}}=\frac{r_E}{\sqrt{c(1-u)+u}},\qquad
 \frac{\dd r_E^{\mathrm{native}}}{\dd r_E}
 =\frac{c}{[c(1-u)+u]^{3/2}}.\label{eq:commonsoft}
\end{equation}
$r_E$ is a unit-interval integration coordinate, not the top-soft factor
$r_t$ introduced later. This coordinate change preserves the real domain
and removes the residual finite logarithm from unequal soft slopes. Native
H-event inversions retain their own recovered coordinates and Born measure.
Selected colour pairs, frame information, sampled coordinates, and soft
conversion factors are saved and restored with the phase-space point.

\section{The combined IF shower density and its singular limits}\label{sec:kernels}

\subsection{One antenna, two FKS orientations}

For the Born quark--quark IF connections, Pythia's local-recoil ISR machinery
describes both incoming and outgoing gluon-emission endpoints. Corresponding
FI contributions in TimeShower are suppressed to avoid duplication. Thus
``ISR and FSR matching'' here means reproducing both physical endpoints of
the combined antenna, not adding two independent shower kernels.

Let $B$ denote the appropriate Born squared matrix element, with its native
flavour and colour-flow information. For the massless-incoming analytic
approximation, the local radiation measure, including the relevant flux
conversion, is
\begin{equation}
 \dd\mu_{\mathrm{IF}}=
 \frac{\dd\zI\,\dd Q^2\,\dd\varphi}{32\pi^3}.
\end{equation}
The raw stock density used for gluon emission is
\begin{equation}
 \boxed{\frac{\mathcal D_0}{B}
 =\frac{2\gs^2}{\zI Q^2}
   \CF\frac{1+\zI^2}{1-\zI},\qquad \CF=\frac43,
 \quad\gs^2=4\pi\as.}\label{eq:stockdensity}
\end{equation}
It is multiplied by the real-state luminosity in the actual counterterm.
The backward incoming rescaling is already represented by the map and
luminosity convention. One must not multiply a second PDF factor $1/\zI$,
the global-recoil shower Jacobian, or a global-recoil \code{xfact} onto
Eq.~\eqref{eq:stockdensity}.

The two FKS orientations divide the same gluon antenna with
\begin{equation}
 \chi_{\mathrm I}=\frac{p_f\cdot k}{p_f\cdot k+p_a\cdot k},
 \qquad
 \chi_{\mathrm F}=\frac{p_a\cdot k}{p_f\cdot k+p_a\cdot k},
 \qquad \chi_{\mathrm I}+\chi_{\mathrm F}=1.
\end{equation}
These partition factors assign a connection to a sector; they do not double
its colour factor. The density summed over the two orientations remains
Eq.~\eqref{eq:stockdensity}.

Backward gluon conversion, $g\to q\bar q$ with a Born incoming quark, uses
\begin{equation}
 \frac{\mathcal D_{g\to q\bar q}}{B}
 =\frac{2\gs^2}{\zI Q^2}\,
 \TR[\zI^2+(1-\zI)^2],\qquad \TR=\frac12.
\end{equation}
It belongs only to the incoming orientation, has no soft-gluon pole, and
requires neither of the local soft/outgoing-collinear G repairs discussed
below. This Born-flavour channel should not be confused with a later shower
state whose current incoming colour end is a gluon.

\subsection{The final-quark collinear numerator}

For $k\parallel p_a$ at fixed $\zI$, $Q^2\to0$ and $\zeta\to1$.
The stock incoming Altarelli--Parisi coefficient is correct.
For a massless outgoing partner, set
\begin{equation}
 s_{fk}=2p_f\cdot k=\frac{wD}{\zI}.
\end{equation}
At the hard final-collinear endpoint $s_{fk}\to0$, $\zI\to1$ while $\zeta$
can remain away from one. The coefficient of $2\gs^2/s_{fk}$ becomes
\begin{align}
 \text{stock:}&\quad \frac{2\CF}{1-\zeta},&
 \text{QCD:}&\quad \CF\frac{1+\zeta^2}{1-\zeta}.\label{eq:finalcol}
\end{align}
The difference is a final-collinear singularity in $\mathcal D_0$, even
though the numerator change looks modest when viewed as a function of two
DIS variables. A matching calculation cannot use bare $R-\mathcal D_0$ and
assume that this difference is integrable.

\subsection{What the paper says about this choice}

Cabouat--Sj\"ostrand, section 3.2, Eq.~(23), keep the usual backward-DGLAP
kernel. They note its two DIS pole locations and the numerator difference
$1+x^2$ versus $x^2+z_1^2$. Their stated motivation for using only the IF
description is that a separate FI description would require a cumbersome
variable transformation and Jacobian, with little benefit for the good
overall description already obtained~\cite{Cabouat}. This is a practical
approximation, not a proof of the exact hard final-quark collinear
coefficient. The proposed numerator patch addresses that specific
coefficient while retaining the common map.

\subsection{The massive-top soft self term}

On the heavy line, the stock soft limit contains the pair eikonal,
\begin{equation}
 \frac{\mathcal D_{0,\mathrm{soft}}}{B}
 =2\gs^2\CF\frac{p_a\cdot p_f}
                    {(p_a\cdot k)(p_f\cdot k)}.
\end{equation}
For $p_f^2=m_t^2$, the required quark-line current is
\begin{equation}
 \boxed{\frac{\mathcal D_{\mathrm{QCD,soft}}}{B}
 =2\gs^2\CF\frac{p_a\cdot p_f}
                    {(p_a\cdot k)(p_f\cdot k)}
  -\gs^2\CF\frac{m_t^2}{(p_f\cdot k)^2}.}\label{eq:softcurrent}
\end{equation}
The second term is the massive quark self-eikonal~\cite{MassiveDipoles}.
The mass in the recoil map and the positivity constraint~\eqref{eq:kt} do
not generate this term automatically in the branching probability.

No massive suppression is applied to the massless light-line shower
density. Exact FKS soft counterterms are sector partitions of the full
Born soft current and can contain information about the massive top even
when inspected in a light-line sector. That is different from dressing a
massless shower line with an ad hoc top-mass factor.

\section{G-function matching with stock Pythia}\label{sec:G}

\subsection{The local effective subtraction}

Let $S$, $C$ and $SC$ denote the exact native FKS soft, collinear and
soft-collinear overlap counterterms. Their definitions include sector
functions, the correct Born projection, and endpoint-to-real measure ratios.
They are not scalar kernels evaluated indiscriminately at one real-state
Born point. For a transition width $h=\code{PythiaLocalGWidth}$, define
\begin{equation}
 g_h(v)=\begin{cases}
 0,&v\leq0,\\
 \displaystyle\frac{u^2}{u^2+(1-u)^2},\quad u=v/h,&0<v<h,\\
 1,&v\geq h.
 \end{cases}\label{eq:Gprofile}
\end{equation}
The default is $h=0.01$. This document uses $g_h$ as the shower multiplier;
the G convention in Eq.~(2.117) of Ref.~\cite{MG5} is complementary.

Suppressing the necessary native measure labels, the implemented stock
prescription is
\begin{equation}
 \boxed{\DG=d(q)\left\{
 g_h(w)\mathcal D_0+
 [1-g_h(w)]\left[S+[1-g_h(1-y_{\mathrm{FKS}})](C-SC)\right]\right\}.}
 \label{eq:DG}
\end{equation}
The angular profile is local to the forward collinear endpoint of each FKS
history. The same effective subtraction is added to S events and subtracted
from H events. The exact overlap subtraction is essential: a pure soft
replacement can itself spoil a collinear limit in its transition region.

The choice $w=1-\zI$ is essential for the stock case. It vanishes both for a
soft gluon and for hard outgoing-collinear radiation. A profile based only on
the real gluon energy would leave the second limit in
Eq.~\eqref{eq:finalcol} unrepaired. The compact profile approaches zero
quadratically, and is exactly one outside the transition region. Gluon
conversion has zero replacement complements.

The starting-scale factor uses $q=\sqrt{wQ^2}$, the shower evolution scale,
not the massive physical $\kt$:
\begin{equation}
 d(q)=1-f\left(\frac{q-\mu_{\min}}{\mu_{\max}-\mu_{\min}}\right),
 \quad
 f(u)=\begin{cases}0&u\leq0,\\
 u^2/[u^2+(1-u)^2]&0<u<1,\\1&u\geq1.
 \end{cases}
\end{equation}
$\mu_{\min},\mu_{\max}$ are the native Born dipole shower-starting-scale
bounds. The raw density and the exact replacement have the same damping.
In the code the returned replacement complement is
$[1-g_h(w)]d(q)$; applying the damping to the raw term alone would leave an
unphysical exact replacement above the shower starting scale.

\subsection{S/H construction and the finite-width accuracy qualification}

For explanation, let $C_{\mathrm{NLO}}$ denote the ordinary NLO subtraction,
$I_C$ its integrated part, and $V$ the renormalized virtual contribution.
Suppressing sums over native maps and their Jacobians, the MC@NLO structure is
\begin{align}
 \overline B_{\mathrm{MC}}
 &=B+V+I_C+\int\dd\Phi_{\mathrm{rad}}
                     (\DG-C_{\mathrm{NLO}}),\\
 H&=R-\DG.
\end{align}
These definitions explain the opposite signs of the MC counterterms and
the cancellation in the inclusive rate. The implementation retains the full
native-history construction of Refs.~\cite{MCatNLO,MG5}; it does not use this
schematic expression to discard history or sector factors.

If the shower actually evolves with $\Dps$ while matching uses $\DG$, the
first nontrivial expansion of an infrared-safe observable contains the
residual term
\begin{equation}
 \delta\sigma[O]=\int\dd\Real\,
  (\Dps-\DG)\,[O(\Real)-O(\Born)]+\cdots.\label{eq:Gpower}
\end{equation}
For stock Pythia this is restricted to the unresolved transition region and
is power suppressed as that region shrinks. At finite $h$ it is not
automatically an extra power of $\as$. A finite-width G repair makes event
weights locally integrable; it does not replace the shower's actual
singular kernels by exact QCD. Width variation is therefore part of the
physics assessment, as emphasized by the general G-function discussion in
Ref.~\cite{MG5}. Reducing $h$ also increases cancellations and can worsen
statistical efficiency.

At the implementation level, complete native contributions are formed
before H-event repartitioning. Native Born amplitudes, luminosities,
sector functions, endpoint measures, and starting scales belong to their
own history. S events use the common soft point of
Eq.~\eqref{eq:commonsoft}. This is why both explicit FKS summation and
explicit MC history summation are required in the local matching mode.
The ordinary local MC@NLO no-emission multiplier is one; the separate Delta
matching machinery is not active.

\section{The two opt-in Pythia patches}\label{sec:patches}

\subsection{Quark IF numerator correction}

The first patch replaces
\begin{equation}
 \CF\frac{1+\zI^2}{w}
 \quad\longrightarrow\quad
 \CF\frac{\zI^2+\zeta^2}{w}.\label{eq:colpatch}
\end{equation}
It is implemented as an additional branching-veto weight
\begin{equation}
 r_q=\frac{\zI^2+\zeta^2}{1+\zI^2}.
\end{equation}
Since $0\leq\zeta\leq1$ on the physical support, $0\leq r_q\leq1$.
The incoming-collinear and soft limits give $r_q\to1$. At the massless
outgoing-collinear limit the numerator becomes $1+\zeta^2$, as required.
The correction is restricted to a current quark $q\to qg$ IF branching
with a final quark or antiquark colour partner. It does not modify II
recoil, incoming-gluon kernels, flavour conversion, masses, or thresholds.

The new setting is
\begin{lstlisting}
SpaceShower:IFquarkCollinearCorrection = on
\end{lstlisting}
and is \emph{off} by default. The corresponding patch file is
\file{Template/NLO/MCatNLO/srcPythia8/patches/pythia8318-local-quark-collinear.patch}.

For the one-patch matching prescription, Eq.~\eqref{eq:DG} uses the corrected
density and replaces $g_h(w)$ by $g_h(\xi_E)$, where
\begin{equation}
 \xi_E=\frac{2P\cdot k}{P^2}
      =\frac{2E_k}{\sqrt{\hat s_{\mathrm R}}},\qquad
 P=p_1+p_2.
\end{equation}
$\xi_E$ is the true energy fraction in the real partonic centre of mass,
independent of which native orientation describes the point. Hard
final-collinear radiation therefore uses the patched density directly.
The top-soft G correction and its $C-SC$ overlap remain necessary.

\subsection{Massive-top soft correction}

The second patch multiplies the corrected heavy-line kernel by
\begin{align}
 \boxed{r_t
 =1-\frac{m_t^2\,p_a\cdot k}
             {2(p_a\cdot p_f)(p_f\cdot k)}
 =1-\frac{m_t^2\zI^2Q^2}{w\zeta D^2}
 =\frac{\kt^2}{wQ^2\zeta}.}\label{eq:rt}
\end{align}
Taking the soft limit of the pair eikonal times $r_t$ gives exactly
Eq.~\eqref{eq:softcurrent}. Physical transverse-momentum support gives
$0\leq r_t\leq1$, and $r_t\to1$ in the incoming-collinear limit.
In the top rest frame the same ratio is $(1+\cos\theta_{ak})/2$, with
$\theta_{ak}$ the angle between the incoming-quark direction and the gluon.
This also makes its positivity transparent.

The added Pythia factor is
\begin{lstlisting}[language=C++]
double z1 = 1. - z * Q2 / (m2ColPair - m2ColPartner);
wt *= pT2corr / (pT2 * z1);
\end{lstlisting}
with $z=\zI$, $z1=\zeta$, $pT2=wQ^2$, and $pT2corr=\kt^2$.
The actual guarded branch requires a current incoming quark with the same
quark flavour before and after the branching, an emitted gluon, and an
outgoing top/antitop colour partner. The new setting is
\begin{lstlisting}
SpaceShower:IFtopSoftCorrection = on
\end{lstlisting}
and is also off by default. The patch file is
\file{Template/NLO/MCatNLO/srcPythia8/patches/pythia8318-local-top-soft.patch}.
It is applied after the numerator patch.

Both factors act inside the branching veto, before an emission is accepted.
Consequently the accepted rate and the no-emission probability are changed
consistently:
\begin{equation}
 \Delta(q_1,q_0)=\exp\left[-\int_{q_0}^{q_1}\dd\Phi_{\mathrm{rad}}
                   K_{\mathrm{modified}}\right].
\end{equation}
Reweighting accepted events alone would not make this Sudakov change.
Because both factors are suppressions, they do not require a larger trial
envelope. Existing stock-Pythia weight-above-unity warnings are a separate
issue and are not hidden by changing that envelope.

With both patches the MG5 raw top-line density is
\begin{equation}
 \frac{\mathcal D_2}{B}=
 \frac{2\gs^2}{\zI Q^2}\,
 \CF\frac{\zI^2+\zeta^2}{w}\,r_t.
\end{equation}
The light line uses the same expression without $r_t$. The required Born
soft and massless collinear limits now agree with QCD within the stated
incoming-mass approximation. Thus the special local G complements are
zero, \code{gfactsf=gfactcl=1}, and \code{PythiaLocalGWidth} is inactive.
Ordinary FKS subtraction and starting-scale damping remain.

\subsection{Subsequent emissions and possible side effects}\label{sec:later}

The top-soft factor follows the \emph{current} colour connection. It is
not limited to the first event emission: radiation on the light line can
leave the incoming-quark--top connection intact. Using that connection's
current momenta keeps the ratio physical when it is encountered later.

After a heavy-line emission creates a final--final gluon--top dipole, this
patch does not act on it. Pythia's existing TimeShower treatment already
includes massive-quark effects and, for a gluon radiating with a massive
quark recoiler, an additional recoil dead-cone suppression~\cite{PythiaTime}.
Multiplying a Born IF factor onto that final--final dipole could double count
mass suppression. \code{SimpleTimeShower.cc} is unchanged, and
\code{TimeShower:recoilDeadCone=on} is retained.

A later incoming-gluon--top IF configuration, for example after backward
flavour conversion, is outside this quark patch's scope. Neither the colour
factor nor the finite splitting function of that configuration should be
inferred from Eq.~\eqref{eq:rt}. No universal repair of every massive shower
topology is claimed.

The bounds on $r_q,r_t$ and the consistent Sudakov establish useful
properties: no negative rate is introduced by these factors, no new
overestimate larger than the original is needed, and the specified Born
singular limits are repaired. They do not prove unchanged logarithmic
accuracy at all orders. Finite hard-emission probabilities and higher-order
shower histories change by design. Matrix-element corrections are kept off;
an existing MEC could otherwise compensate, conflict with, or duplicate
the deliberately modified base rate. Phenomenological effects of the two
patches should be assessed with the matched samples, not dismissed as
automatically negligible.

\subsection{Isolation, input parameters, and unmodified effects}

The private source changes touch \code{SimpleSpaceShower.cc} and the XML
settings declaration only. They do not change recoil kinematics or quark
masses. A private build with the flags off reproduces the reference shower
at the checked points. The installed prefixes are
\begin{lstlisting}
stock:       /tmp/singletop_local_matching_pythia
one patch:   HEPTools/pythia8_local_collinear_patch
two patches: HEPTools/pythia8_local_2patch
\end{lstlisting}
The shared stock installation is never patched and rebuilt around a run.
The selected library search path is local to the shower subprocess. Thus
success, failure, interruption, and concurrent runs do not require a
shared-library restoration step. Source and library hashes record this.

Several separate effects remain distinct from the kernel repairs:
\begin{itemize}
\item Shower masses and heavy-flavour thresholds differ from a massless
five-flavour incoming calculation. These familiar mass/threshold effects
are retained in Pythia; no blanket zero-mass override is used.
\item Smooth ISR screening, shower cutoffs, couplings, PDFs, MEC switches,
and recoil options are steering choices. The chosen settings are documented
below. Changing a numerator or adding the massive self term is a source-level
kernel change, not just another existing input value.
\item The extra ISR screening is disabled for this comparison. The ordinary
shower termination cutoff remains; the analytic subtraction stays
unregularized in the unresolved region.
\item Pythia's azimuthal/coherence models are not promoted to exact
fixed-order spin correlations by these scalar kernel patches. Small effects
ordinarily neglected in MC@NLO remain under that qualification.
\end{itemize}

\section{Implementation structure and runtime contract}\label{sec:implementation}

\subsection{Maps, native states, and subtraction}

The main implementation responsibilities are as follows; all paths are
relative to the repository root.
\begin{description}[style=nextline]
\item[\file{Template/NLO/SubProcesses/FKSParams.f90}]
Declares, reads, prints and validates the local-map and paired-patch flags.
The collinear patch requires local recoil; the soft patch requires the
collinear patch. Local matching requires the explicit sums and Pythia8.
\item[\file{Template/NLO/SubProcesses/mc_native_runtime.f}]
\code{mc_select_pythia_pair} selects the unique IF partner from immutable
native Born IDs and colour metadata. It supports the two beam orderings,
antiquarks and crossed real channels, including an incoming gluon with a
quark Born projection. It does not guess a recoil partner from a real-state
particle number alone.
\item[\file{Template/NLO/SubProcesses/genps_fks_radiation.f}]
\code{generate_momenta_pythia_initial} and
\code{invert_momenta_pythia_initial} implement ISR generation and inversion.
\code{pythia_local_radiation_coordinates} supplies the common S-event soft
point and its Jacobian. Final-state orientations reuse the incoming-reservoir
FSR construction, including massive solutions and frame conversions.
\item[\file{Template/NLO/SubProcesses/recoil_selection.f}]
Selects the local policy and rejects combinations with explicit recoil
requests or resonance recoil.
\item[\file{Template/NLO/SubProcesses/fks_phase_space_data.f}]
Carries the selected pair, local ISR coordinate and soft conversion with the
complete point state. Born reuse checks the chosen pair as well as the
usual kinematic context.
\item[\file{Template/NLO/SubProcesses/fks_singular.f}]
Implements the real and endpoint prefactors, including the finite soft
conversion and the native FKS contributions used in the G replacement.
\item[\file{Template/NLO/SubProcesses/montecarlocounter.f}]
\file{prepare_pythia_local_kinematics} computes stable invariants;
\file{compute_pythia_local_counterterm} selects the stock, one-patch or
two-patch kernel and the relevant G profile. It keeps the selected Born
weights, orientation partition, luminosity convention and scale damping
together.
\end{description}
Supporting native-history, event-colour and starting-scale code carries the
same Born context through S and H contributions. The changes were kept in
the phase-space and matching owners rather than made as a correction to
an already generated event file.

Unsupported Born structures fail explicitly: the present selector expects
two massless incoming quarks in the Born channel, one massive outgoing top,
one massless outgoing quark, and one colour flow containing the two distinct
IF connections. The process need not have the same real-channel incoming
flavours as its Born projection. Hadron beam and QCD-only requirements are
also checked.

\subsection{FKS flags and launcher checks}

For local matching, the common settings in \file{Cards/FKS_params.dat} are
\begin{lstlisting}
#FKSPythiaLocalRecoil
1
#FKSISRMapping
0
#FKSFinalRecoil
0
#FKSExplicitSum
.true.
#MCExplicitKLSum
.true.
#PythiaLocalGWidth
0.01d0
\end{lstlisting}
The one-patch process additionally sets
\code{PythiaLocalCollinearPatch=.true.}; the two-patch process also sets
\code{PythiaLocalSoftPatch=.true.}. An absent flag in an older process card
uses its compiled default, false. The default-recoil process does not turn
on the local mode and keeps its ordinary summation settings.

The shower launcher is
\file{Template/NLO/MCatNLO/Scripts/MCatNLO_MadFKS_PYTHIA8.Script}.
It reads the paired FKS flags, checks that the selected XML declares the
required shower setting, and writes the matching flag into
\file{Pythia8.cmd}. It writes a declared optional flag off when the paired
matching prescription is off, and avoids passing an unknown flag to stock
Pythia. The selected build is given by \code{pythia8_path} in
\file{Cards/amcatnlo_configuration.txt}. The process-local shower template
puts the selected private library ahead of inherited library paths.

Set \code{parton_shower=PYTHIA8}, \code{ickkw=0}, and
\code{mcatnlo_delta=False}. Separate initial-only/final-only matching jobs
are not a substitute for the combined IF prescription. Do not combine the
local mode with nonzero generic ISR/final recoil overrides. Regenerate
integration grids and LHE events after changing the map, kernel-patch flags,
or G width. Showering an old LHE sample with a newly patched kernel is not
the corresponding matched prediction.

\section{The common production and shower setup}\label{sec:setup}

The current four run cards and parameter cards are byte-identical.
They request $10^6$ events each at $\sqrt s=13$ TeV with $m_t=173\GeV$,
\code{pdlabel=nn23nlo}, and dynamic scales
\begin{equation}
 \mu_R=\mu_F=H_T/2,
 \qquad H_T=\sum_{i\in\mathrm{final}}
                  \sqrt{m_i^2+p_{\perp i}^2}.
\end{equation}
The card value is \code{dynamical_scale_choice=-1}; this scale definition is
given by the generated \code{setscales.f}. Scale/PDF reweighting and merging
are off. The listed \code{lhaid} is not an active LHAPDF request when the
built-in \code{nn23nlo} label is selected.

\textbf{Current versus historical cuts.} The current cards set
\code{ptj=-1} and \code{etaj=-1}: the generation jet cuts are disabled.
The card's generation jet algorithm/radius are not the analysis definition.
Older fixed-order and G-width validation examples used \code{ptj=10} GeV.
Their cross sections are reported below as historical checks, not directly
compared with the current unrestricted-generation rates.

The common shower choices are
\begin{center}\small
\begin{tabularx}{\linewidth}{@{}lX@{}}
\toprule
Item & Value / interpretation\\
\midrule
QCD ISR and FSR & On\\
Hadronisation and hadron-level processing & Off\\
MPI, remnants, primordial transverse momentum & Off\\
Colour reconnection & Off\\
QED showers from quarks, leptons, photons and other particles & Off\\
Weak showers and top decays & Off; \code{6:mayDecay=off}\\
ISR/FSR matrix-element corrections & Off, including extended / after-first options\\
Shower $\as(M_Z)$ & 0.118, one-loop running, CMW rescaling off\\
Shower PDFs & Internal \code{PDF:pSet=15}, NNPDF2.3 NLO QCD+QED grid\\
Extra ISR screening & \code{pT0parametrization=0}, \code{pT0Ref force=0.}\\
FSR gluon-to-quark weight & \code{TimeShower:weightGluonToQuark=1}\\
FSR recoil dead cone & On\\
QED conversion upper mass & \code{TimeShower:mMaxGamma=10}; inactive when QED showers are off\\
Shower card events / jobs & All available events, one shower job, seed 172925\\
\bottomrule
\end{tabularx}
\end{center}
The PDF grid's fitted coupling and the chosen shower coupling are distinct
inputs; the latter is the standard MG5 shower value used consistently in
the comparison. No LHAPDF installation is needed for the internal shower
grid.

Both showers use \code{pTmaxMatch=1} and unit \code{pTmaxFudge}, with
\code{shower_scale_factor=1.0} in the run card. The card's SCALUP/start-scale
prescription is retained. The common options also include
\code{SpaceShower:pTdampMatch=0}, \code{SpaceShower:samePTasMPI=off}, and
\code{TimeShower:allowBeamRecoil=on}. Event normalization is
\code{event_norm=average}; \code{nsplit_jobs=1} and \code{nevents=-1} in the
shower card request the whole generated sample in one job.

The recoil switches distinguish the samples:
\begin{lstlisting}
# pp_tj_default
SpaceShower:dipoleRecoil = off
TimeShower:globalRecoil = on

# all three local samples
SpaceShower:dipoleRecoil = on
TimeShower:globalRecoil = off
\end{lstlisting}
The default sample retains the usual MG5 first-emission global-recoil
steering, including its S/H-related options. This is not a statement that
every subsequent FSR branching uses global recoil. Pythia's ordinary
later shower machinery remains in use.

The final event is an open partonic subsystem because beam remnants are
off. Its full-proton charge and momentum check is therefore disabled with
\code{Check:event=off}. This is a specific event-record choice; successful
acceptance in these production runs is not a full-event momentum-check
validation. Closed-record tests with remnants on are discussed separately.

The extra screening assignment uses \code{DM_1} so that Pythia's
case-sensitive lower-case \code{force} survives MG5's handling of the
\code{pythia8_options} dictionary. \code{DM_2} supplies an absolute XML path,
which also makes the sibling PDF-data directory available from a shower run
directory. If installations move, update the Pythia prefix, this XML/data
path, and the HepMC include/library paths together.

\section{Analysis definitions and the purpose of every plot family}\label{sec:analysis}

\subsection{Objects and selections}

The common analysis is
\file{Template/NLO/MCatNLO/PY8Analyzer/py8an_HwU_pp_tj_recoil.f}, copied to
each process. It runs through the normal MG5--Pythia8 interface and fills
HwU histograms. Final-state quarks, including bottom, and gluons are
clustered with anti-$k_T$, $R=0.4$, using FastJet or the bundled fjcore.
The stable top is excluded. The analysis requires exactly one final stable
top and rejects final hadrons, photons, or leptons.

Accepted jets satisfy $p_{\perp j}>20\GeV$ and $|y_j|<4.5$.
$j_1,j_2,\ldots$ are the accepted jets ordered by transverse momentum.
$j_1$ is a kinematic proxy for the spectator jet, not an ancestry or flavour
tag; hard radiation can become $j_1$. No infrared-unsafe light-parton flavour
assignment is used.

There are three selections:
\begin{align}
 \mathrm{ALL}:&\quad\text{no accepted-jet requirement},\\
 \mathrm{TAG}:&\quad p_{\perp j_1}>30\GeV,\\
 \mathrm{GAP}:&\quad\mathrm{TAG},\quad
            |y_t-y_{j_1}|>3,\quad m(t,j_1)>500\GeV.
\end{align}
These are analysis cuts, not generation cuts. A plot requiring $j_2$ is
filled only when at least two accepted jets are present. Events failing
TAG or GAP still count towards the inclusive event normalization.

True rapidity and pseudorapidity are
\begin{equation}
 y(p)=\frac12\log\frac{E+p_z}{E-p_z},\qquad
 \eta(p)=\operatorname{arsinh}\frac{p_z}{p_\perp}.
\end{equation}
The top is massive, and clustered jets can also have nonzero mass; $y$ and
$\eta$ should not be interchanged in the gap definitions.

\subsection{Standard single-top distributions}

The ALL selection contains the total rate, top transverse momentum and
rapidity, $|y_t|$, a wider top-$p_\perp$ spectrum, jet multiplicity, jet
$H_T$, top-plus-jets $H_T$, and leading-jet $p_\perp$ when a jet exists.
TAG and GAP contain their fiducial rates, top and leading-jet $p_\perp/y$,
leading-jet $\eta$ and $|\eta|$, and the following standard correlation
observables:
\begin{align}
 \Delta y_{t1}&=y_t-y_{j_1},&
 \Delta\phi_{t1}&=\arccos\frac{\boldsymbol p_{\perp t}\cdot
                  \boldsymbol p_{\perp j_1}}{p_{\perp t}p_{\perp j_1}},\\
 m_{t1}&=\sqrt{(p_t+p_{j_1})^2},&
 y_{t1}&=y(p_t+p_{j_1}).
\end{align}
The azimuthal difference lies in $[0,\pi]$. These plots test inclusive
production and migration across fiducial cuts, rather than only a special
recoil-sensitive corner. Multiplicity histograms put all $N_{\rm jets}\geq8$
in the last bin.

\subsection{Recoil and additional-jet observables}

At Born level the top and its light recoil are balanced. The spectra of
\begin{equation}
 p_\perp(t+j_1),\qquad
 p_\perp(t+j_1+j_2)
\end{equation}
measure recoil from radiation not included in those reconstructed systems.
The first is also booked with fine bins below $50\GeV$. The second requires
$j_2$ and probes additional shower radiation and radiation outside accepted
jets. The low-imbalance region is sensitive to Sudakov evolution, jet
clustering and recoil; a hard tail can also reflect which jet is leading.

The analogy to Higgs VBF concerns the \emph{first additional jet}: $j_2$
here corresponds to VBF's $j_3$ when the two VBF tagging jets are treated
as the Born visible system. VBF studies motivate rapidity centrality and
gap activity as probes of coherence and recoil~\cite{VBFJager,VBFHoche}.
The massive coloured top makes the present process different; neither the
size nor the sign of an effect is inferred from the analogy alone.

Define
\begin{equation}
 Y=|y_{j_1}-y_t|,\quad
 y_m=\frac{y_t+y_{j_1}}2,\quad
 s=\operatorname{sign}(y_{j_1}-y_t).
\end{equation}
The analysis fills $j_2$ $p_\perp,y,\eta$, the unscaled displacement
$y_{j_2}-y_m$, and the oriented Zeppenfeld coordinate
\begin{equation}
 z^*_{j_2}=s\frac{y_{j_2}-y_m}{Y}.\label{eq:zstar}
\end{equation}
The top and leading jet are at $-1/2$ and $+1/2$, respectively, independent
of beam reversal. $|z^*|<1/2$ puts $j_2$ in the rapidity interval between
them; $z^*=0$ is its midpoint. The $z^*$ histograms require $Y>1$ to avoid
unstable small denominators. A second version requires
$p_{\perp j_2}>30\GeV$. These conditions are in addition to TAG or GAP.

The oriented distances
\begin{equation}
 s(y_{j_2}-y_{j_1}),\qquad s(y_{j_2}-y_t)
\end{equation}
separate activity near the light and heavy ends without averaging them away
under beam reversal. The signed distance to the nearest gap edge is
\begin{equation}
 d_{\rm edge}=\begin{cases}
 +\min(|y_{j_2}-y_t|,|y_{j_2}-y_{j_1}|),&j_2\text{ inside the gap},\\
 -\min(|y_{j_2}-y_t|,|y_{j_2}-y_{j_1}|),&j_2\text{ outside the gap}.
 \end{cases}
\end{equation}
This distinguishes central in-gap radiation from radiation just beyond a
tagging object without division by $Y$.

\subsection{Gap activity and veto efficiency}

Extra jets are $j_2,j_3,\ldots$. The open interval between $y_t$ and
$y_{j_1}$ defines gap membership. The analysis measures the number of gap
jets, their scalar transverse-momentum sum, their hardest transverse
momentum, and extra-jet $H_T$ in $|y_j-y_m|<0.5$. Events with no such activity
contribute at zero. The last gap multiplicity bin includes $N_{\rm gap}\geq5$.

For a veto threshold $Q_0$, define
\begin{equation}
 \epsilon_{\rm gap}(Q_0)
 =\frac{\sigma(\text{selection and no extra gap jet with }p_\perp\geq Q_0)}
        {\sigma(\text{selection})}.
\end{equation}
The code tests the hardest extra gap jet with a strict $p_\perp<Q_0$
condition and fills cumulative numerator and reference scans at
$Q_0=21,23,\ldots,99\GeV$. The normal plots turn these four raw scan
histograms into two efficiencies, one for TAG and one for GAP. They are
not differential cross sections in the threshold.

The complete catalogue in Appendix~\ref{app:catalogue} lists all binning and
ranges. There are nine ALL histograms and 33 each for TAG and GAP: 75 raw
histograms. The normal process PDF has 37 pages because identical titles
overlay selections. The four-sample comparison has 73 panels:
$75-4+2$, replacing the four raw veto scans with two efficiencies.

\subsection{Weights, normalization, and uncertainty propagation}\label{sec:statistics}

Signed MC@NLO weights are retained. The central contribution uses the
HepMC event weight; variations are normalized relatively to the central
LHE weight. The analysis calls \code{HwU_add_points} once per analysed event,
including cut failures. It counts actual analysis calls for normalization,
not the event serial number. The HwU error mode is explicitly Poisson:
the squared signed weights are accumulated. Thus a negative weight reduces
the bin content while increasing its variance.

For a numerator $N$ that is a weighted subset of a denominator $D$, the
Poisson convention gives $\operatorname{Cov}(N,D)=\operatorname{Var}(N)$.
The comparison helper propagates
\begin{equation}
 f=\frac ND,\qquad
 (\delta f)^2=\frac{(1-f)^2\operatorname{Var}(N)
              +f^2[\operatorname{Var}(D)-\operatorname{Var}(N)]}{D^2}.
\end{equation}
The variance difference is floored at zero against roundoff. This is used
for gap efficiencies and for bin fractions in shape-normalized spectra.
Between separately generated samples, ratios use the independent-sample
formula
\begin{equation}
 R=\frac AB,\qquad
 \delta R=\sqrt{\frac{(\delta A)^2}{B^2}
                   +\frac{A^2(\delta B)^2}{B^4}}.
\end{equation}
Nonpositive denominators are masked. A tiny positive denominator can still
produce an enormous, meaningless ratio in a sparsely populated or strongly
cancelling tail. Some automatically scaled atlas ratio panels consequently
have very large ranges. Those tails should be rebinned or restricted before
quantitative interpretation; the displayed ratio scale is not evidence of
a huge physical effect.

Absolute HwU contents are bin-integrated cross sections. Comparison plots
divide by bin width for a differential density. With \code{--shape}, each
distribution is normalized over its displayed range, including the subset
correlation above. One-bin rates and gap efficiencies are not shape
normalized. Histogram statistical errors are not scale/PDF uncertainties,
integration errors, or a theory-error estimate.

The existing Pythia driver checks its selected-event limit before analysing
the last selected event. Thus the 1000-, 10000- and million-event runs
normally contain 999, 9999 and 999999 analysis calls. The analysis
normalizes with its own call count. This pre-existing one-event effect was
documented; it is not the HEPEVT truncation failure.

\section{Current production results and how to read them}\label{sec:results}

\textbf{Correction of 25 October 2026.}
The MC@NLO samples in this section and the comparison atlas must be
regenerated. The final clustering vertex stored its daughter momenta but
omitted the mother's momentum. For a coloured core mother, as in
$Wb\to t$, the two-coloured-leg scale
$q^2=2|p_b\cdot p_t|$ therefore used uninitialized data. Different S and H
evaluations could then use different shower damping scales, preventing
their MC subtraction terms from cancelling. Initializing
\file{p_inter(:,0,0)=pcl(:,3)} in \file{cluster.f} fixes the defect.
Compiled tests cover the $Wb\to t$ core and $bu\to td$ clustering with
both beam orientations; independent integrations verify inclusive
subtraction closure. The prepared MC@NLO sources and executables have
been updated. Saved events and plots are preserved as historical
diagnostics, not validated predictions. Fixed-order and POWHEG results
are unaffected. Numerical checks are recorded under
\file{docs/audits/singletop_normalization_mcexplicit.json}.

The following values are taken from the four current
\file{Events/run_01/plot_PYTHIA8_1_0.HwU} files, as summarized in
\file{Analyses/pp_tj_recoil.json}. Each corresponding shower log reports
one million accepted events. These are the superseded production samples,
subsequent to the short validation runs described later.

\begin{table}[htbp]
\centering\small
\begin{tabular}{@{}lrrrr@{}}
\toprule
Observable & Default & Local & 1 patch & 2 patches\\
\midrule
ALL: total & $138.452\pm0.342$ & $138.424\pm0.291$ & $137.508\pm0.254$ & $138.283\pm0.233$ \\
TAG: fiducial & $104.592\pm0.285$ & $102.857\pm0.233$ & $102.198\pm0.203$ & $102.777\pm0.184$ \\
GAP: fiducial & $19.101\pm0.127$ & $18.806\pm0.104$ & $18.698\pm0.091$ & $18.748\pm0.084$ \\
TAG: at least two jets & $51.649\pm0.196$ & $49.152\pm0.146$ & $48.646\pm0.126$ & $47.222\pm0.116$ \\
GAP: at least two jets & $8.111\pm0.084$ & $7.205\pm0.056$ & $7.028\pm0.048$ & $6.774\pm0.045$ \\
\bottomrule
\end{tabular}
\caption{Current shower-analysis rates in pb. Errors are Poisson histogram statistical errors. All four production samples contain one million selected shower events.}\label{tab:production}
\end{table}


The inclusive predictions are close on the percent scale, while the
additional-jet selections show larger differences. The rates must not be
interpreted as a precision estimate of a recoil systematic: the samples
have different matched maps/counterterms, the stock local case has a finite
G width, and the patched cases deliberately change the shower. The quoted
errors in this table are histogram sampling errors only.

Equal generated event counts do not give equal signed-weight precision.
For the inclusive Poisson histogram, define
\begin{equation}
 N_{\rm eff}=\frac{(\sum_i w_i)^2}{\sum_i w_i^2}
            =\left(\frac{\sigma_{\rm ALL}}
                         {\delta\sigma_{\rm ALL}}\right)^2.
\end{equation}
The current inclusive bins give approximately 164000, 226000, 292000 and
353000 effective events for default, local, one patch and two patches,
respectively. These are measures of signed-weight sampling precision, not
counts of positive events; the effective count depends on the observable.

\begin{table}[htbp]
\centering\small
\begin{tabular}{@{}lrrrr@{}}
\toprule
Observable & Default & Local & 1 patch & 2 patches\\
\midrule
TAG: central j2 & $10.094\pm0.103$ & $8.168\pm0.059$ & $8.012\pm0.050$ & $7.910\pm0.046$ \\
GAP: central j2 & $3.426\pm0.060$ & $2.459\pm0.032$ & $2.378\pm0.027$ & $2.374\pm0.025$ \\
\bottomrule
\end{tabular}
\caption{Cross section in pb for $|z^*_{j_2}|<1/2$, from the displayed centrality histogram. TAG also requires $Y>1$; GAP already has $Y>3$. At least two accepted jets are required.}\label{tab:central}
\end{table}


In the GAP selection the central-$j_2$ rate is about $28$--$31\%$ lower in
the three local predictions than in the default prediction. This is an
observed difference in these samples, not a deduction from the VBF papers.
The extra top-soft patch leaves this particular central-$j_2$ rate close
to the one-patch value while changing the total two-jet rate. A full shape
comparison is therefore more informative than one number.

\begin{table}[htbp]
\centering\small
\begin{tabular}{@{}lrrrr@{}}
\toprule
Observable & Default & Local & 1 patch & 2 patches\\
\midrule
TAG: Q0=21 GeV & $0.8565\pm0.0010$ & $0.8897\pm0.0006$ & $0.8923\pm0.0005$ & $0.8940\pm0.0005$ \\
TAG: Q0=41 GeV & $0.9443\pm0.0006$ & $0.9564\pm0.0004$ & $0.9570\pm0.0003$ & $0.9575\pm0.0003$ \\
TAG: Q0=61 GeV & $0.9745\pm0.0004$ & $0.9807\pm0.0003$ & $0.9805\pm0.0002$ & $0.9807\pm0.0002$ \\
GAP: Q0=21 GeV & $0.7890\pm0.0029$ & $0.8564\pm0.0017$ & $0.8614\pm0.0014$ & $0.8610\pm0.0013$ \\
GAP: Q0=41 GeV & $0.9251\pm0.0017$ & $0.9413\pm0.0011$ & $0.9434\pm0.0009$ & $0.9447\pm0.0008$ \\
GAP: Q0=61 GeV & $0.9680\pm0.0010$ & $0.9759\pm0.0007$ & $0.9757\pm0.0006$ & $0.9761\pm0.0006$ \\
\bottomrule
\end{tabular}
\caption{Gap-veto efficiencies at three thresholds, with the subset covariance included.}\label{tab:veto}
\end{table}


At $Q_0=21\GeV$, the GAP veto efficiency rises from about $0.789$ in the
default prediction to $0.856$--$0.861$ in the local predictions. At larger
thresholds the difference decreases. This is consistent with reduced
accepted in-gap activity for these local setups, and complements the
rapidity-centrality plots.

The following reading order is useful.
\begin{enumerate}
\item Check the total, TAG and GAP rates, top $p_\perp/y$, leading-jet
$p_\perp/\eta$, and $m(t,j_1)$. These establish the usual single-top
production and acceptance changes.
\item Compare $|\Delta y_{t1}|$, the $j_2$ displacement from the midpoint,
and $z^*_{j_2}$, especially in GAP. The signed coordinates distinguish
activity near the two different colour lines.
\item Use gap multiplicity, hardest gap-jet $p_\perp$, gap/midpoint $H_T$
and veto efficiencies to determine whether a centrality change comes with
more or less gap radiation.
\item Compare absolute and shape-normalized spectra. A lower $j_2$ rate
and a changed normalized centrality distribution are separate effects.
\item Inspect pair/system transverse momentum and $\Delta\phi$ for recoil
balance, with appropriate caution in low-statistics tails.
\end{enumerate}

Top-plus-leading-jet observables have the usual nominal NLO+PS status,
subject to the explicit finite-width and shower-model qualifications above.
One additional resolved jet is supplied by the NLO real term and shower,
and has lower fixed-order accuracy. More additional jets are shower
dominated. There is no multijet merging. Neither the atlas nor these
central values establishes a preferred shower without uncertainty studies.

\section{Validation evidence and its limits}\label{sec:validation}

\subsection{Fixed-order map validation}

The local-map tests compare forward/inverse closure, preserved masses and
spectators, allowed support, Bjorken fractions and independent Cartesian
phase-space determinants. They include both beam orders, antiquarks,
emitted-leg permutations, native replay, state restoration, massive
solutions and invalid-configuration rejection. Invalid arithmetic,
division by zero and overflow are trapped in the relevant numerical tests.
The finite soft conversion is checked against the dimensionally regulated
energy integral, and the incoming-collinear/PDF prefactors against the
ordinary ISR map.

The earlier independent-seed fixed-order production check used the same
13 TeV process but the earlier validation cuts, including $p_{\perp j}>10$
GeV at generation. It gave
\begin{center}
\begin{tabular}{@{}lrrr@{}}
\toprule
Map & Seed & $(\xi_S^{\rm cut},\delta_S)$ & NLO cross section [pb]\\
\midrule
Default & 7251 & $(0.5,1)$ & $130.7\pm0.46$\\
Local & 8313 & $(0.5,1)$ & $129.0\pm0.55$\\
Local, varied subtraction cutoffs & 9301 & $(0.25,0.5)$ & $129.9\pm0.70$\\
\bottomrule
\end{tabular}
\end{center}
The default/local difference is 2.37 combined standard errors; the local
cutoff variation is 1.01. The differential comparison had 406 tested bin
differences across the two comparisons, with maximum marginal pulls 2.92
and 3.11, below the stated 3.84 Bonferroni threshold. These are scoped
numerical consistency checks, not a proof from inclusive agreement alone.

\subsection{Stock-G matching and width variation}

All four generated subprocess groups pass the production ME and MC
soft/collinear tests and 20/20 pole checks. Additional checks compare the
soft-collinear overlap in common coordinates. One apparent soft offset
was traced to diagnostic limit flags being cleared during inversion and
not restored for amplitude evaluation: one side used a $10^{-6}$ shortcut
while the reference used the diagnostic $10^{-12}$ threshold. That
diagnostic inconsistency was corrected; it was not cured with a physics
factor.

The historical stock-G study with the earlier generation cuts found
\begin{center}
\begin{tabular}{@{}rrrr@{}}
\toprule
$h$ & Integration [pb] & Showered events & Negative-weight events\\
\midrule
0.005 & $132.1\pm2.5$ & 1000 & 247\\
0.010 & $134.7\pm1.6$ & 1000 & 267\\
0.020 & $134.9\pm1.7$ & 1000 & 238\\
\bottomrule
\end{tabular}
\end{center}
All LHE weights and scales were finite. The widths were statistically
compatible at that limited precision; this does not measure a precise
differential G-width uncertainty. Stock Pythia completed the samples but
reported retries, momentum-check, threshold, dipole-mass and occasional
weight-above-unity messages. G functions do not repair shower evolution
or trial-envelope problems.

The stock audit covers 1200 accepted first emissions across both beams and
quarks/antiquarks, and compares 100 accepted gluon-emission points with the
Fortran raw density. It establishes the checked map, routing and analytic
massless-incoming density. It is not a full audit of every finite
acceptance factor or every heavy-flavour threshold branch.

\subsection{Patch and complete-history checks}

The one-patch audit compares 100 accepted patched points with the Fortran
density and 100 disabled-flag points bit for bit against stock. The
two-patch audit verifies corresponding opt-out equivalence with both stock
and the one-patch build, 100 two-patch density points, invariant bounds,
Lorentz invariance, and the intended flavour/topology restrictions.
Complete histories include later quark--top IF branchings and final--final
gluon--top emissions, retaining the ordinary recoil dead-cone setting.
The generated one- and two-patch 1000-event validation runs both pass the
four subprocess ME/MC checks and pole checks, shower all 1000 events, and
write all 75 histograms.

The two-patch kernel/launcher test batch had 31 passing tests and its
map/matching regression batch 46. The subsequent record-fix batch had 15
passing tests, including the long-history regression. These are recorded
batches at development checkpoints, with overlapping tests; the counts
should not be added as independent evidence. The earlier stock-G audit
also records a broader 155-test batch.

Closed proton-record diagnostics enable beam remnants, keep primordial
$k_\perp$ off, and turn \code{Check:event} on. Their event-check rejections
must be reported as failures of that check:
\begin{center}
\begin{tabular}{@{}lrrr@{}}
\toprule
Fixture & Stock & Collinear patch & Both patches\\
\midrule
Four beam/flavour cases, 1200 attempts per variant & 21 & 33 & 22\\
One case, 1000 attempts per variant & 10 & 19 & 20\\
\bottomrule
\end{tabular}
\end{center}
No top-soft-factor bound or stable-top-mass violations were found in those
checks. This is not equivalent to every complete event passing momentum
conservation. Failures also occur in the references; these finite samples
neither prove equal rejection probabilities nor identify a new failure
specific to the top-soft factor. The production open-record setup has
different checks and cannot settle this question by its acceptance count.

\section{The default-shower failure and the HEPEVT fix}\label{sec:hepevt}

The first full default shower stopped at event 2866 with
\code{Expected exactly one stable final top}. Replaying that event with
the original settings gave
\begin{center}
\begin{tabular}{@{}lrrr@{}}
\toprule
 & Pythia & HepMC & Fortran HEPEVT\\
\midrule
Stable final tops & 1 & 1 & 0\\
History entries & 5174 & 5173 & 4000\\
\bottomrule
\end{tabular}
\end{center}
Pythia's entry zero accounts for its extra entry. Global recoil stores many
intermediate spectator copies, so the full history can be much longer than
the final partonic state. HepMC's \code{IO_HEPEVT} silently truncated the
history at the configured 4000 entries and lost the final top. Pythia had
not decayed or removed it.

The stable-top analysis now uses
\file{Template/NLO/MCatNLO/include/HEPMC_ST_RECOIL.INC} with 100000 entries.
At \code{PYABEG} it calls \code{hepmc_set_max_number_entries}, defined in
\file{Template/NLO/MCatNLO/include/HEPEVTSize.h}, so HepMC uses the same
array strides as the Fortran common block. Increasing only one side would
misalign the record. The helper validates the declared capacity against
the allocation supported by the HepMC headers. The analysis also refuses
a record at the capacity limit, instead of using a possibly truncated state.

This analysis and interface fix are installed in all four process
directories and the export template. Ordinary other analyses keep the
generic 4000-entry layout unless they opt in. HepMC-output-only drivers do
not use this Fortran analysis transfer. Pythia source/library kernels,
analysis cuts and original LHE events are unchanged.

The first 10000 original events were replayed successfully. Event 2866 now
retains all 5173 HepMC entries and one final top in each representation;
the largest replayed record has 5177 entries. The usual MG5 shower command
also accepted all 10000 events and produced 75 histograms and a 37-page
PDF. A regression places the stable top beyond history entry 4000 and
checks unchanged histogram results. The later completed default production
shower reports one million accepted events; its result is the default
curve in the current atlas. Earlier audit text saying the million-event
rerun had not yet happened describes the earlier fix checkpoint.

\section{Running, rebuilding, reproducing, and extending the study}\label{sec:running}

From the default or stock-local process directory, shower existing events
with
\begin{lstlisting}
./bin/shower run_01 -f
\end{lstlisting}
From either patched process, a fresh matched generation and automatic
shower are requested by
\begin{lstlisting}
./bin/generate_events -f
\end{lstlisting}
The first available run name is chosen. For existing events in a patched
process, \code{./bin/shower run_01 -f} also selects its private shower.
Outputs have the form
\begin{lstlisting}
Events/run_01/plot_PYTHIA8_<shower-index>_0.HwU
Events/run_01/plot_PYTHIA8_<shower-index>_0.pdf
Events/run_01/plot_PYTHIA8_<shower-index>_0.html
\end{lstlisting}
Plotting scripts for Matplotlib and GnuPlot are also saved. The comparison
helper accepts process directories or explicit HwU files and selects the
latest shower by modification time when given a process directory.

From the repository root:
\begin{lstlisting}
python3 Analyses/compare_pp_tj_recoil.py
python3 Analyses/compare_pp_tj_recoil.py --patched
python3 Analyses/compare_pp_tj_recoil.py --two-patches
python3 Analyses/compare_pp_tj_recoil.py --two-patches --shape
\end{lstlisting}
The last option includes all four samples. Without \code{--shape}, the
default outputs are \file{Analyses/pp_tj_recoil.pdf} and its JSON rate table;
shape outputs have the suffix \code{_shapes}. The standard colours are blue
for default, orange for local stock, green for the collinear patch and purple
for both patches. Ratios are always to default. \code{--run} and
\code{--out} select another run or output prefix. Explicit HwU files are
useful when short validation runs have different names in each process.

Private builds can be reproduced in \emph{new} prefixes outside the stock
source tree:
\begin{lstlisting}
python3 docs/audits/build_pythia8_local_collinear.py STOCK NEW_ONE_PATCH
python3 docs/audits/build_pythia8_local_collinear.py STOCK NEW_TWO_PATCHES --top-soft
\end{lstlisting}
The builder verifies reference 8.318 source hashes, copies files without
hard-linking mutable sources, applies patches in order, rebuilds the static
and shared libraries, and records commands and SHA256 values in
\code{local_collinear_patch.json}. Its final stock-hash check also runs on
failure. A differently patched or differently versioned reference is not
silently accepted.

The main focused checks can be run with
\begin{lstlisting}
python3 -m unittest \
 tests.unit_tests.fks.test_pythia_local_recoil \
 tests.unit_tests.fks.test_pythia_local_matching \
 tests.unit_tests.fks.test_recoil_selection \
 tests.unit_tests.various.test_pythia8 \
 tests.unit_tests.various.test_single_top_recoil_analysis -v

python3 docs/audits/check_pythia8_singletop_local.py STOCK
python3 docs/audits/check_pythia8_local_collinear_patch.py STOCK ONE_PATCH
python3 docs/audits/check_pythia8_local_soft_patch.py STOCK ONE_PATCH TWO_PATCHES
\end{lstlisting}
The C++/Fortran checks require the appropriate compilers and installed
Pythia/HepMC dependencies. The read-only stock audit can additionally take
an LHE file and steering file for complete-shower checks.

The most useful next physics studies are systematic G-width variations
for the stock local and one-patch prescriptions, shower starting-scale and
coupling variations, robustly rebinned shape comparisons, and a dedicated
analysis of later massive gluon--top IF configurations. A separate study
would be needed for top decays, resonance matching, antitop phenomenology,
four-flavour production, merging, or general colour flows. None is implied
by the present stable-top validation.

To rebuild this report from the repository root, use
\begin{lstlisting}
cd docs
pdflatex -interaction=nonstopmode -halt-on-error single_top_local_recoil.tex
pdflatex -interaction=nonstopmode -halt-on-error single_top_local_recoil.tex
\end{lstlisting}
The source contains the prose, equations, numerical tables, bibliography,
histogram catalogue and plot index. It imports the existing comparison PDF
for the atlas; the compiled report embeds those pages and is self-contained
for reading. Recompiling against changed plot files without updating the
numerical snapshot tables would make the report inconsistent; regenerate
or review those tables together with the atlas.

\begin{thebibliography}{9}
\addcontentsline{toc}{section}{References}
\bibitem{Cabouat}
B. Cabouat and T. Sj\"ostrand,
\emph{Some Dipole Shower Studies}, Eur. Phys. J. C 78 (2018) 226,
\href{https://arxiv.org/abs/1710.00391}{arXiv:1710.00391};
especially sections 3.1--3.4 and Eq.~(23).
\bibitem{FKS}
S. Frixione, Z. Kunszt and A. Signer,
\emph{Three-jet cross sections to next-to-leading order},
Nucl. Phys. B467 (1996) 399--442,
\href{https://arxiv.org/abs/hep-ph/9512328}{hep-ph/9512328}.
\bibitem{MCatNLO}
S. Frixione and B. R. Webber,
\emph{Matching NLO QCD computations and parton shower simulations},
JHEP 06 (2002) 029,
\href{https://arxiv.org/abs/hep-ph/0204244}{hep-ph/0204244}.
\bibitem{MG5}
J. Alwall et al.,
\emph{The automated computation of tree-level and next-to-leading order
differential cross sections, and their matching to parton shower simulations},
JHEP 07 (2014) 079,
\href{https://arxiv.org/abs/1405.0301}{arXiv:1405.0301};
especially Eq.~(2.117) and the G-function discussion on pp.~48--49.
\bibitem{MassiveDipoles}
S. Catani, S. Dittmaier, M. H. Seymour and Z. Tr\'ocs\'anyi,
\emph{The Dipole Formalism for Next-to-Leading Order QCD Calculations with
Massive Partons}, Nucl. Phys. B627 (2002) 189--265,
\href{https://arxiv.org/abs/hep-ph/0201036}{hep-ph/0201036}.
\bibitem{PythiaTime}
Pythia manual, \emph{Timelike Showers}, sections on global recoil and
\code{TimeShower:recoilDeadCone},
\url{https://pythia.org/latest-manual/TimelikeShowers.html}.
The implementation used here is specifically Pythia 8.318, whose local
source/XML is the versioned reference.
\bibitem{VBFJager}
B. J\"ager, A. Karlberg, S. Pl\"atzer, J. Scheller and M. Zaro,
\emph{Parton-shower effects in Higgs production via Vector-Boson Fusion},
Eur. Phys. J. C 80 (2020) 756,
\href{https://arxiv.org/abs/2003.12435}{arXiv:2003.12435}.
\bibitem{VBFHoche}
S. H\"oche, S. Mrenna, S. Payne, C. T. Preuss and P. Skands,
\emph{A Study of QCD Radiation in VBF Higgs Production with Vincia and Pythia},
\href{https://arxiv.org/abs/2106.10987}{arXiv:2106.10987}.
\end{thebibliography}

\clearpage\appendix
\section{Source, test, and data provenance}\label{app:provenance}

This report consolidates the implementation and evidence in the following
records; their essential content is reproduced in the body, so they need
not be read to understand the report.
\begin{itemize}
\item \file{docs/native_fks_projection.md}: map, measure, finite endpoint,
native replay, and the earlier fixed-order numerical comparisons.
\item \file{docs/audits/pythia8318_singletop_local.md} and
\file{docs/audits/pythia8318_singletop_local_validation.json}: stock-G
prescription, current checks, width samples, and explicitly historical
withdrawn-shower results.
\item \file{docs/audits/pythia8318_local_collinear_patch.txt} and
\file{docs/audits/pythia8318_singletop_local_patch_validation.json}:
one-patch construction and validation.
\item \file{docs/audits/pythia8318_local_2patch.txt} and
\file{docs/audits/pythia8318_singletop_local_2patch_validation.json}:
top-soft factor, complete-history counts and two-patch checks.
\item \file{docs/audits/pp_tj_default_hepevt_fix.txt} and
\file{docs/audits/pp_tj_default_hepevt_validation.json}: the record-transfer
failure and its 10000-event regression checkpoint.
\item \file{Analyses/README.single_top_recoil},
\file{Analyses/compare_pp_tj_recoil.py}, and the common Fortran analysis:
plot definitions, error treatment and comparison commands.
\end{itemize}

The principal numerical fixture sources are
\begin{lstlisting}
tests/input_files/check_pythia_local_recoil.f90
tests/input_files/check_pythia_local_matching.f90
tests/input_files/check_pythia8_singletop_local.cc
tests/input_files/check_pythia8_local_soft_history.cc
tests/input_files/check_pythia8_singletop_lhe.cc
tests/unit_tests/fks/test_pythia_local_recoil.py
tests/unit_tests/fks/test_pythia_local_matching.py
tests/unit_tests/various/test_pythia8.py
tests/unit_tests/various/test_single_top_recoil_analysis.py
\end{lstlisting}
The production plot inputs and their fingerprints at the document snapshot
are listed below. All four process directories use the same run and
parameter card contents; their FKS prescriptions and Pythia build selections
differ as described in Table~\ref{tab:variants}.

\begin{lstlisting}
pp_tj_default/Events/run_01/plot_PYTHIA8_1_0.HwU
SHA256 b6285ce6ec3f79074fdbf009f5029d3f8d802d78f84ff383a08b79da1ca66cdc
pp_tj_local/Events/run_01/plot_PYTHIA8_1_0.HwU
SHA256 d8ab24388950604abb4d49d6c44fe53fdb7c4891b1d377ce2f65e33e14db5058
pp_tj_local_patch/Events/run_01/plot_PYTHIA8_1_0.HwU
SHA256 a250298c45d0075560d626ee394bb2e882f2dee044c838879f3b8dd2e6087992
pp_tj_local_2patch/Events/run_01/plot_PYTHIA8_1_0.HwU
SHA256 8ae580ebe6fbe130d70fb57b695f90260543599e7633db92fda8551614826061
Analyses/pp_tj_recoil.pdf
SHA256 d44ee8022815b89e48e019eb10f264b86ed9f5335f7978f96e3ea52ddb3984bc
common Cards/run_card.dat
SHA256 4a99dde5381a2313afbcc588b0df4d289a3287cd4c6dc573cb88ebe2614a07cc
common Cards/param_card.dat
SHA256 bcc9d7f3d1e249ec6cc2e6de24e7511c62da576a7812ae17cf492ce7d3f0f41f
\end{lstlisting}


\clearpage
\section{Complete histogram catalogue}\label{app:catalogue}

The labels below are the exact HwU titles, with the selection suffix omitted.
Each ALL row is one histogram. Every row in the TAG/GAP table is booked
twice, once for each selection. Bounds are histogram edges; out-of-range
entries do not define an extrapolation beyond the displayed spectrum.
The multiplicity plots explicitly clamp their overflow categories to
the last booked bin. The upper azimuth edge is $\pi+10^{-10}$ in the code
to include the endpoint safely.

\subsection*{ALL: nine histograms}
\begin{longtable}{@{}L{88mm}rrr@{}}
\toprule
Title & Bins & Lower edge & Upper edge\\
\midrule
\endfirsthead
\toprule
Title & Bins & Lower edge & Upper edge\\
\midrule
\endhead
\bottomrule
\endfoot
total rate [pb] & 1 & 0 & 1 \\
top pT [GeV] & 60 & 0 & 600 \\
top y & 50 & -5 & 5 \\
top abs y & 25 & 0 & 5 \\
Njets20 (last >=8) & 9 & -0.5 & 8.5 \\
jets HT [GeV] & 50 & 0 & 1000 \\
top plus jets HT [GeV] & 50 & 0 & 1200 \\
top pT high [GeV] & 50 & 0 & 1500 \\
j1 pT [GeV] & 60 & 0 & 600 \\
\end{longtable}
\subsection*{TAG and GAP: 33 histograms per selection}
\begin{longtable}{@{}L{88mm}rrr@{}}
\toprule
Title & Bins & Lower edge & Upper edge\\
\midrule
\endfirsthead
\toprule
Title & Bins & Lower edge & Upper edge\\
\midrule
\endhead
\bottomrule
\endfoot
top pT [GeV] & 60 & 0 & 600 \\
top y & 50 & -5 & 5 \\
top abs y & 25 & 0 & 5 \\
Njets20 (last >=8) & 9 & -0.5 & 8.5 \\
j1 pT [GeV] & 60 & 0 & 600 \\
fiducial rate [pb] & 1 & 0 & 1 \\
j1 y & 50 & -5 & 5 \\
j1 eta & 50 & -5 & 5 \\
j1 abs eta & 25 & 0 & 5 \\
abs dy(t,j1) & 40 & 0 & 10 \\
dy(t,j1) & 50 & -10 & 10 \\
dphi(t,j1) [rad] & 32 & 0 & $\pi+10^{-10}$ \\
pT(t+j1) [GeV] & 60 & 0 & 300 \\
pT(t+j1) low [GeV] & 50 & 0 & 50 \\
m(t,j1) [GeV] & 60 & 150 & 1650 \\
y(t+j1) & 50 & -5 & 5 \\
j2 pT [GeV] & 50 & 0 & 300 \\
j2 y & 50 & -5 & 5 \\
j2 eta & 50 & -5 & 5 \\
j2 y minus midpoint & 50 & -5 & 5 \\
j2 zstar (dy>1) & 40 & -2 & 2 \\
j2 zstar pT>30 (dy>1) & 40 & -2 & 2 \\
j2 signed gap distance & 50 & -5 & 5 \\
j2 dy from j1 (oriented) & 50 & -6 & 6 \\
j2 dy from top (oriented) & 50 & -6 & 6 \\
pT(t+j1+j2) [GeV] & 60 & 0 & 300 \\
Ngap jets20 (last >=5) & 6 & -0.5 & 5.5 \\
hardest gap jet pT [GeV] & 50 & 0 & 200 \\
gap jets HT [GeV] & 50 & 0 & 300 \\
midpoint extra jets HT [GeV] & 50 & 0 & 300 \\
gap veto survived Q0 [GeV] & 40 & 20 & 100 \\
gap veto reference Q0 [GeV] & 40 & 20 & 100 \\
at least two jets rate [pb] & 1 & 0 & 1 \\
\end{longtable}


\clearpage
\section{Comparison atlas: index and complete figures}\label{app:atlas}

Every panel below compares all four million-event matched predictions.
The upper panel shows an absolute cross section or density (or a veto
efficiency); the lower panel shows ratios to default. Colours are blue,
orange, green and purple in the order of Table~\ref{tab:variants}.
Error bars are the statistical errors described in
section~\ref{sec:statistics}. No scale/PDF band is included. Very small
positive denominators can dominate the automatic ratio-axis range in
tails; consult the absolute distribution and statistical errors before
interpreting them. The index provides links to each embedded panel.

\begin{longtable}{@{}rL{112mm}r@{}}
\toprule
Panel & Selection and observable & Report page\\
\midrule
\endfirsthead
\toprule
Panel & Selection and observable & Report page\\
\midrule
\endhead
\bottomrule
\endfoot
1 & \hyperref[atlas:p01]{ALL: total rate [pb]} & \pageref{atlas:p01}\\
2 & \hyperref[atlas:p02]{ALL: top pT [GeV]} & \pageref{atlas:p02}\\
3 & \hyperref[atlas:p03]{TAG: top pT [GeV]} & \pageref{atlas:p03}\\
4 & \hyperref[atlas:p04]{GAP: top pT [GeV]} & \pageref{atlas:p04}\\
5 & \hyperref[atlas:p05]{ALL: top y} & \pageref{atlas:p05}\\
6 & \hyperref[atlas:p06]{TAG: top y} & \pageref{atlas:p06}\\
7 & \hyperref[atlas:p07]{GAP: top y} & \pageref{atlas:p07}\\
8 & \hyperref[atlas:p08]{ALL: top abs y} & \pageref{atlas:p08}\\
9 & \hyperref[atlas:p09]{TAG: top abs y} & \pageref{atlas:p09}\\
10 & \hyperref[atlas:p10]{GAP: top abs y} & \pageref{atlas:p10}\\
11 & \hyperref[atlas:p11]{ALL: Njets20 (last >=8)} & \pageref{atlas:p11}\\
12 & \hyperref[atlas:p12]{TAG: Njets20 (last >=8)} & \pageref{atlas:p12}\\
13 & \hyperref[atlas:p13]{GAP: Njets20 (last >=8)} & \pageref{atlas:p13}\\
14 & \hyperref[atlas:p14]{ALL: jets HT [GeV]} & \pageref{atlas:p14}\\
15 & \hyperref[atlas:p15]{ALL: top plus jets HT [GeV]} & \pageref{atlas:p15}\\
16 & \hyperref[atlas:p16]{ALL: top pT high [GeV]} & \pageref{atlas:p16}\\
17 & \hyperref[atlas:p17]{ALL: j1 pT [GeV]} & \pageref{atlas:p17}\\
18 & \hyperref[atlas:p18]{TAG: j1 pT [GeV]} & \pageref{atlas:p18}\\
19 & \hyperref[atlas:p19]{GAP: j1 pT [GeV]} & \pageref{atlas:p19}\\
20 & \hyperref[atlas:p20]{TAG: fiducial rate [pb]} & \pageref{atlas:p20}\\
21 & \hyperref[atlas:p21]{GAP: fiducial rate [pb]} & \pageref{atlas:p21}\\
22 & \hyperref[atlas:p22]{TAG: j1 y} & \pageref{atlas:p22}\\
23 & \hyperref[atlas:p23]{GAP: j1 y} & \pageref{atlas:p23}\\
24 & \hyperref[atlas:p24]{TAG: j1 eta} & \pageref{atlas:p24}\\
25 & \hyperref[atlas:p25]{GAP: j1 eta} & \pageref{atlas:p25}\\
26 & \hyperref[atlas:p26]{TAG: j1 abs eta} & \pageref{atlas:p26}\\
27 & \hyperref[atlas:p27]{GAP: j1 abs eta} & \pageref{atlas:p27}\\
28 & \hyperref[atlas:p28]{TAG: abs dy(t,j1)} & \pageref{atlas:p28}\\
29 & \hyperref[atlas:p29]{GAP: abs dy(t,j1)} & \pageref{atlas:p29}\\
30 & \hyperref[atlas:p30]{TAG: dy(t,j1)} & \pageref{atlas:p30}\\
31 & \hyperref[atlas:p31]{GAP: dy(t,j1)} & \pageref{atlas:p31}\\
32 & \hyperref[atlas:p32]{TAG: dphi(t,j1) [rad]} & \pageref{atlas:p32}\\
33 & \hyperref[atlas:p33]{GAP: dphi(t,j1) [rad]} & \pageref{atlas:p33}\\
34 & \hyperref[atlas:p34]{TAG: pT(t+j1) [GeV]} & \pageref{atlas:p34}\\
35 & \hyperref[atlas:p35]{GAP: pT(t+j1) [GeV]} & \pageref{atlas:p35}\\
36 & \hyperref[atlas:p36]{TAG: pT(t+j1) low [GeV]} & \pageref{atlas:p36}\\
37 & \hyperref[atlas:p37]{GAP: pT(t+j1) low [GeV]} & \pageref{atlas:p37}\\
38 & \hyperref[atlas:p38]{TAG: m(t,j1) [GeV]} & \pageref{atlas:p38}\\
39 & \hyperref[atlas:p39]{GAP: m(t,j1) [GeV]} & \pageref{atlas:p39}\\
40 & \hyperref[atlas:p40]{TAG: y(t+j1)} & \pageref{atlas:p40}\\
41 & \hyperref[atlas:p41]{GAP: y(t+j1)} & \pageref{atlas:p41}\\
42 & \hyperref[atlas:p42]{TAG: j2 pT [GeV]} & \pageref{atlas:p42}\\
43 & \hyperref[atlas:p43]{GAP: j2 pT [GeV]} & \pageref{atlas:p43}\\
44 & \hyperref[atlas:p44]{TAG: j2 y} & \pageref{atlas:p44}\\
45 & \hyperref[atlas:p45]{GAP: j2 y} & \pageref{atlas:p45}\\
46 & \hyperref[atlas:p46]{TAG: j2 eta} & \pageref{atlas:p46}\\
47 & \hyperref[atlas:p47]{GAP: j2 eta} & \pageref{atlas:p47}\\
48 & \hyperref[atlas:p48]{TAG: j2 y minus midpoint} & \pageref{atlas:p48}\\
49 & \hyperref[atlas:p49]{GAP: j2 y minus midpoint} & \pageref{atlas:p49}\\
50 & \hyperref[atlas:p50]{TAG: j2 zstar (dy>1)} & \pageref{atlas:p50}\\
51 & \hyperref[atlas:p51]{GAP: j2 zstar (dy>1)} & \pageref{atlas:p51}\\
52 & \hyperref[atlas:p52]{TAG: j2 zstar pT>30 (dy>1)} & \pageref{atlas:p52}\\
53 & \hyperref[atlas:p53]{GAP: j2 zstar pT>30 (dy>1)} & \pageref{atlas:p53}\\
54 & \hyperref[atlas:p54]{TAG: j2 signed gap distance} & \pageref{atlas:p54}\\
55 & \hyperref[atlas:p55]{GAP: j2 signed gap distance} & \pageref{atlas:p55}\\
56 & \hyperref[atlas:p56]{TAG: j2 dy from j1 (oriented)} & \pageref{atlas:p56}\\
57 & \hyperref[atlas:p57]{GAP: j2 dy from j1 (oriented)} & \pageref{atlas:p57}\\
58 & \hyperref[atlas:p58]{TAG: j2 dy from top (oriented)} & \pageref{atlas:p58}\\
59 & \hyperref[atlas:p59]{GAP: j2 dy from top (oriented)} & \pageref{atlas:p59}\\
60 & \hyperref[atlas:p60]{TAG: pT(t+j1+j2) [GeV]} & \pageref{atlas:p60}\\
61 & \hyperref[atlas:p61]{GAP: pT(t+j1+j2) [GeV]} & \pageref{atlas:p61}\\
62 & \hyperref[atlas:p62]{TAG: Ngap jets20 (last >=5)} & \pageref{atlas:p62}\\
63 & \hyperref[atlas:p63]{GAP: Ngap jets20 (last >=5)} & \pageref{atlas:p63}\\
64 & \hyperref[atlas:p64]{TAG: hardest gap jet pT [GeV]} & \pageref{atlas:p64}\\
65 & \hyperref[atlas:p65]{GAP: hardest gap jet pT [GeV]} & \pageref{atlas:p65}\\
66 & \hyperref[atlas:p66]{TAG: gap jets HT [GeV]} & \pageref{atlas:p66}\\
67 & \hyperref[atlas:p67]{GAP: gap jets HT [GeV]} & \pageref{atlas:p67}\\
68 & \hyperref[atlas:p68]{TAG: midpoint extra jets HT [GeV]} & \pageref{atlas:p68}\\
69 & \hyperref[atlas:p69]{GAP: midpoint extra jets HT [GeV]} & \pageref{atlas:p69}\\
70 & \hyperref[atlas:p70]{TAG: at least two jets rate [pb]} & \pageref{atlas:p70}\\
71 & \hyperref[atlas:p71]{GAP: at least two jets rate [pb]} & \pageref{atlas:p71}\\
72 & \hyperref[atlas:p72]{TAG: gap veto efficiency} & \pageref{atlas:p72}\\
73 & \hyperref[atlas:p73]{GAP: gap veto efficiency} & \pageref{atlas:p73}\\
\end{longtable}


\clearpage
\includepdf[pages={1},pagecommand={\thispagestyle{plain}\phantomsection\label{atlas:p01}}]{../Analyses/pp_tj_recoil.pdf}
\includepdf[pages={2},pagecommand={\thispagestyle{plain}\phantomsection\label{atlas:p02}}]{../Analyses/pp_tj_recoil.pdf}
\includepdf[pages={3},pagecommand={\thispagestyle{plain}\phantomsection\label{atlas:p03}}]{../Analyses/pp_tj_recoil.pdf}
\includepdf[pages={4},pagecommand={\thispagestyle{plain}\phantomsection\label{atlas:p04}}]{../Analyses/pp_tj_recoil.pdf}
\includepdf[pages={5},pagecommand={\thispagestyle{plain}\phantomsection\label{atlas:p05}}]{../Analyses/pp_tj_recoil.pdf}
\includepdf[pages={6},pagecommand={\thispagestyle{plain}\phantomsection\label{atlas:p06}}]{../Analyses/pp_tj_recoil.pdf}
\includepdf[pages={7},pagecommand={\thispagestyle{plain}\phantomsection\label{atlas:p07}}]{../Analyses/pp_tj_recoil.pdf}
\includepdf[pages={8},pagecommand={\thispagestyle{plain}\phantomsection\label{atlas:p08}}]{../Analyses/pp_tj_recoil.pdf}
\includepdf[pages={9},pagecommand={\thispagestyle{plain}\phantomsection\label{atlas:p09}}]{../Analyses/pp_tj_recoil.pdf}
\includepdf[pages={10},pagecommand={\thispagestyle{plain}\phantomsection\label{atlas:p10}}]{../Analyses/pp_tj_recoil.pdf}
\includepdf[pages={11},pagecommand={\thispagestyle{plain}\phantomsection\label{atlas:p11}}]{../Analyses/pp_tj_recoil.pdf}
\includepdf[pages={12},pagecommand={\thispagestyle{plain}\phantomsection\label{atlas:p12}}]{../Analyses/pp_tj_recoil.pdf}
\includepdf[pages={13},pagecommand={\thispagestyle{plain}\phantomsection\label{atlas:p13}}]{../Analyses/pp_tj_recoil.pdf}
\includepdf[pages={14},pagecommand={\thispagestyle{plain}\phantomsection\label{atlas:p14}}]{../Analyses/pp_tj_recoil.pdf}
\includepdf[pages={15},pagecommand={\thispagestyle{plain}\phantomsection\label{atlas:p15}}]{../Analyses/pp_tj_recoil.pdf}
\includepdf[pages={16},pagecommand={\thispagestyle{plain}\phantomsection\label{atlas:p16}}]{../Analyses/pp_tj_recoil.pdf}
\includepdf[pages={17},pagecommand={\thispagestyle{plain}\phantomsection\label{atlas:p17}}]{../Analyses/pp_tj_recoil.pdf}
\includepdf[pages={18},pagecommand={\thispagestyle{plain}\phantomsection\label{atlas:p18}}]{../Analyses/pp_tj_recoil.pdf}
\includepdf[pages={19},pagecommand={\thispagestyle{plain}\phantomsection\label{atlas:p19}}]{../Analyses/pp_tj_recoil.pdf}
\includepdf[pages={20},pagecommand={\thispagestyle{plain}\phantomsection\label{atlas:p20}}]{../Analyses/pp_tj_recoil.pdf}
\includepdf[pages={21},pagecommand={\thispagestyle{plain}\phantomsection\label{atlas:p21}}]{../Analyses/pp_tj_recoil.pdf}
\includepdf[pages={22},pagecommand={\thispagestyle{plain}\phantomsection\label{atlas:p22}}]{../Analyses/pp_tj_recoil.pdf}
\includepdf[pages={23},pagecommand={\thispagestyle{plain}\phantomsection\label{atlas:p23}}]{../Analyses/pp_tj_recoil.pdf}
\includepdf[pages={24},pagecommand={\thispagestyle{plain}\phantomsection\label{atlas:p24}}]{../Analyses/pp_tj_recoil.pdf}
\includepdf[pages={25},pagecommand={\thispagestyle{plain}\phantomsection\label{atlas:p25}}]{../Analyses/pp_tj_recoil.pdf}
\includepdf[pages={26},pagecommand={\thispagestyle{plain}\phantomsection\label{atlas:p26}}]{../Analyses/pp_tj_recoil.pdf}
\includepdf[pages={27},pagecommand={\thispagestyle{plain}\phantomsection\label{atlas:p27}}]{../Analyses/pp_tj_recoil.pdf}
\includepdf[pages={28},pagecommand={\thispagestyle{plain}\phantomsection\label{atlas:p28}}]{../Analyses/pp_tj_recoil.pdf}
\includepdf[pages={29},pagecommand={\thispagestyle{plain}\phantomsection\label{atlas:p29}}]{../Analyses/pp_tj_recoil.pdf}
\includepdf[pages={30},pagecommand={\thispagestyle{plain}\phantomsection\label{atlas:p30}}]{../Analyses/pp_tj_recoil.pdf}
\includepdf[pages={31},pagecommand={\thispagestyle{plain}\phantomsection\label{atlas:p31}}]{../Analyses/pp_tj_recoil.pdf}
\includepdf[pages={32},pagecommand={\thispagestyle{plain}\phantomsection\label{atlas:p32}}]{../Analyses/pp_tj_recoil.pdf}
\includepdf[pages={33},pagecommand={\thispagestyle{plain}\phantomsection\label{atlas:p33}}]{../Analyses/pp_tj_recoil.pdf}
\includepdf[pages={34},pagecommand={\thispagestyle{plain}\phantomsection\label{atlas:p34}}]{../Analyses/pp_tj_recoil.pdf}
\includepdf[pages={35},pagecommand={\thispagestyle{plain}\phantomsection\label{atlas:p35}}]{../Analyses/pp_tj_recoil.pdf}
\includepdf[pages={36},pagecommand={\thispagestyle{plain}\phantomsection\label{atlas:p36}}]{../Analyses/pp_tj_recoil.pdf}
\includepdf[pages={37},pagecommand={\thispagestyle{plain}\phantomsection\label{atlas:p37}}]{../Analyses/pp_tj_recoil.pdf}
\includepdf[pages={38},pagecommand={\thispagestyle{plain}\phantomsection\label{atlas:p38}}]{../Analyses/pp_tj_recoil.pdf}
\includepdf[pages={39},pagecommand={\thispagestyle{plain}\phantomsection\label{atlas:p39}}]{../Analyses/pp_tj_recoil.pdf}
\includepdf[pages={40},pagecommand={\thispagestyle{plain}\phantomsection\label{atlas:p40}}]{../Analyses/pp_tj_recoil.pdf}
\includepdf[pages={41},pagecommand={\thispagestyle{plain}\phantomsection\label{atlas:p41}}]{../Analyses/pp_tj_recoil.pdf}
\includepdf[pages={42},pagecommand={\thispagestyle{plain}\phantomsection\label{atlas:p42}}]{../Analyses/pp_tj_recoil.pdf}
\includepdf[pages={43},pagecommand={\thispagestyle{plain}\phantomsection\label{atlas:p43}}]{../Analyses/pp_tj_recoil.pdf}
\includepdf[pages={44},pagecommand={\thispagestyle{plain}\phantomsection\label{atlas:p44}}]{../Analyses/pp_tj_recoil.pdf}
\includepdf[pages={45},pagecommand={\thispagestyle{plain}\phantomsection\label{atlas:p45}}]{../Analyses/pp_tj_recoil.pdf}
\includepdf[pages={46},pagecommand={\thispagestyle{plain}\phantomsection\label{atlas:p46}}]{../Analyses/pp_tj_recoil.pdf}
\includepdf[pages={47},pagecommand={\thispagestyle{plain}\phantomsection\label{atlas:p47}}]{../Analyses/pp_tj_recoil.pdf}
\includepdf[pages={48},pagecommand={\thispagestyle{plain}\phantomsection\label{atlas:p48}}]{../Analyses/pp_tj_recoil.pdf}
\includepdf[pages={49},pagecommand={\thispagestyle{plain}\phantomsection\label{atlas:p49}}]{../Analyses/pp_tj_recoil.pdf}
\includepdf[pages={50},pagecommand={\thispagestyle{plain}\phantomsection\label{atlas:p50}}]{../Analyses/pp_tj_recoil.pdf}
\includepdf[pages={51},pagecommand={\thispagestyle{plain}\phantomsection\label{atlas:p51}}]{../Analyses/pp_tj_recoil.pdf}
\includepdf[pages={52},pagecommand={\thispagestyle{plain}\phantomsection\label{atlas:p52}}]{../Analyses/pp_tj_recoil.pdf}
\includepdf[pages={53},pagecommand={\thispagestyle{plain}\phantomsection\label{atlas:p53}}]{../Analyses/pp_tj_recoil.pdf}
\includepdf[pages={54},pagecommand={\thispagestyle{plain}\phantomsection\label{atlas:p54}}]{../Analyses/pp_tj_recoil.pdf}
\includepdf[pages={55},pagecommand={\thispagestyle{plain}\phantomsection\label{atlas:p55}}]{../Analyses/pp_tj_recoil.pdf}
\includepdf[pages={56},pagecommand={\thispagestyle{plain}\phantomsection\label{atlas:p56}}]{../Analyses/pp_tj_recoil.pdf}
\includepdf[pages={57},pagecommand={\thispagestyle{plain}\phantomsection\label{atlas:p57}}]{../Analyses/pp_tj_recoil.pdf}
\includepdf[pages={58},pagecommand={\thispagestyle{plain}\phantomsection\label{atlas:p58}}]{../Analyses/pp_tj_recoil.pdf}
\includepdf[pages={59},pagecommand={\thispagestyle{plain}\phantomsection\label{atlas:p59}}]{../Analyses/pp_tj_recoil.pdf}
\includepdf[pages={60},pagecommand={\thispagestyle{plain}\phantomsection\label{atlas:p60}}]{../Analyses/pp_tj_recoil.pdf}
\includepdf[pages={61},pagecommand={\thispagestyle{plain}\phantomsection\label{atlas:p61}}]{../Analyses/pp_tj_recoil.pdf}
\includepdf[pages={62},pagecommand={\thispagestyle{plain}\phantomsection\label{atlas:p62}}]{../Analyses/pp_tj_recoil.pdf}
\includepdf[pages={63},pagecommand={\thispagestyle{plain}\phantomsection\label{atlas:p63}}]{../Analyses/pp_tj_recoil.pdf}
\includepdf[pages={64},pagecommand={\thispagestyle{plain}\phantomsection\label{atlas:p64}}]{../Analyses/pp_tj_recoil.pdf}
\includepdf[pages={65},pagecommand={\thispagestyle{plain}\phantomsection\label{atlas:p65}}]{../Analyses/pp_tj_recoil.pdf}
\includepdf[pages={66},pagecommand={\thispagestyle{plain}\phantomsection\label{atlas:p66}}]{../Analyses/pp_tj_recoil.pdf}
\includepdf[pages={67},pagecommand={\thispagestyle{plain}\phantomsection\label{atlas:p67}}]{../Analyses/pp_tj_recoil.pdf}
\includepdf[pages={68},pagecommand={\thispagestyle{plain}\phantomsection\label{atlas:p68}}]{../Analyses/pp_tj_recoil.pdf}
\includepdf[pages={69},pagecommand={\thispagestyle{plain}\phantomsection\label{atlas:p69}}]{../Analyses/pp_tj_recoil.pdf}
\includepdf[pages={70},pagecommand={\thispagestyle{plain}\phantomsection\label{atlas:p70}}]{../Analyses/pp_tj_recoil.pdf}
\includepdf[pages={71},pagecommand={\thispagestyle{plain}\phantomsection\label{atlas:p71}}]{../Analyses/pp_tj_recoil.pdf}
\includepdf[pages={72},pagecommand={\thispagestyle{plain}\phantomsection\label{atlas:p72}}]{../Analyses/pp_tj_recoil.pdf}
\includepdf[pages={73},pagecommand={\thispagestyle{plain}\phantomsection\label{atlas:p73}}]{../Analyses/pp_tj_recoil.pdf}


\end{document}
